\documentclass{article}
\usepackage[left=2cm, right=2cm, top=2cm, bottom=2cm]{geometry}
\usepackage{graphicx}
\usepackage{amsmath}
\usepackage{amssymb}
\usepackage{amsthm}
\usepackage{fancyhdr}
\usepackage{verbatim}
\usepackage{listings}
\usepackage{xcolor}
\usepackage{pdfpages}
\usepackage{cancel}
\usepackage{physics}

\lstset{
    language=C++,
    basicstyle=\ttfamily\footnotesize,
    keywordstyle=\color{blue},
    commentstyle=\color{green},
    stringstyle=\color{red},
    numbers=left,
    numberstyle=\tiny\color{gray},
    frame=single,
    breaklines=true,
    captionpos=b,
}


\title{MTH 553: Midterm \# 1 Recap}
\author{Cliff Sun}

\newtheorem{theorem}{Theorem}[section]
\newtheorem{lemma}[theorem]{Lemma}
\newtheorem{definition}[theorem]{Definition}
\newtheorem{conjecture}[theorem]{Conjecture}
\newtheorem{proposition}[theorem]{Proposition}
\newtheorem{corollary}[theorem]{Corollary}
\newtheorem{one minute paper}[theorem]{One Minute Paper}

\pagestyle{fancy}
\lhead{\textbf{\thepage}\ \ \nouppercase{\rightmark}}
\chead{MTH 553: Midterm \# 1 Recap}
\rhead{Cliff Sun}

\begin{document}

\maketitle

\section*{Wave equation}

Homogenous wave equation on an interval with 
\begin{enumerate}
    \item Dirichlet boundary conditions (Fourier Series with sin)
    \item Neumann boundary conditions (Fourier Series with cos)
\end{enumerate}

How does a wave reflect with 
\begin{enumerate}
    \item Dirichlet BC 
    \item Neumann BC
\end{enumerate}

\section*{Weak derivatives and weak solutions}

Transmission condition: jump must remain constant along the discontinuity. For a wave equation, the jump must remain constant on the characteristics ($x \pm ct$). 

\section*{Higher dimension wave equation}

Spherical mean is the average of the function over a sphere of radius $r$, called $M_h(x,r)$ is the average at point $x$ with radius $r$. Then, used Darboux to prove that $$\triangle M_h = M_{\triangle h}$$

Then, $M_h(x,r)/r$ satisified the wave equation for one dimension where $x$ is fixed. Then Kirchhof showed that the domain of dependence and the range of influence was the 
surface of the cone which implies sharp signals in three dimensions.

For 2 dimension, get the Hadamard formula by dimension reduction. Then the domain of dependence and the range of influence are \textbf{solid} cones, therefore, there are non sharp signals. In n dimensions, 
we also had \textbf{energy methods}. That is 

\begin{equation}
    0 = \frac{1}{2}\frac{d}{dt}\int_{\mathbb{R}}\left( u_t^2 + c^2 |\nabla u|^2 \right)dx
\end{equation}

Prove that energy is conserved. Dispersion $\omega(k)$, plug in the plane wave 

\begin{equation}
    e^{i(kx - \omega t)}
\end{equation}

Moreover, dissipation means that the energy is decreasing. 

\section*{First order equations}

Method of characteristics. 

\begin{equation}
    au_x + bu_y = c
\end{equation}

Special cases: (1) burger's equation and (2) traffic flow.

\begin{equation}
    uu_x + u_y = 0, \text{ Burger's equation}
\end{equation}

Here, $G(u) = u^2/2$. 

\begin{equation}
    G(\rho)_x + \rho_y = 0, \text{ Traffic equation}
\end{equation}

Where 

\begin{equation}
    G(\rho) = c\left( 1- \frac{\rho}{\rho_{\max}} \right)
\end{equation}

All examples of conservation laws. The characteristics are all straight lines, that is 

\begin{equation}
    x = G'(u)y + x_0
\end{equation}

Problems, shocks can occur and also fans. For shocks, we care about 

\begin{enumerate}
    \item $\xi'(y)$
    \item Entropy condition
\end{enumerate}

For a fan, the formula 

\begin{equation}
    G'^{-1}(x/y)
\end{equation}

describes the fan from the origin by Riemann's problem. Also figure out the bounds! 

\section*{ODE methods}

\begin{enumerate}
    \item Separable
    \item 1st order linear
\end{enumerate}

Bounds, for the fan, we have the characteristic is $x = G'(u) + x_0$. Then $x < G'(u_l)y$ is where the left region is, $G'(u_l)y < x < G'(u_r)y$, and $x > G'(u_r)$. 

\end{document}